\documentclass{ximera}
\title{Product and quotient rules}
\begin{document}
\begin{abstract}
Why $(fg)' \neq f'g'$, the product rule and the quotient rule, with an application to Ohm's law. Textbook §3.2.
\end{abstract}
\maketitle

Recall that when we did addition and subtraction of derivatives, we saw that a very similar thing happened as with the limit laws; we could just separate them. For the limit laws we had:
\[
\lim _{x \rightarrow a}(f(x)+g(x))=\left(\lim _{x \rightarrow a} f(x)\right)+\left(\lim _{x \rightarrow a} g(x)\right)
\]

And for derivatives we had:
\[
\frac{d}{d x}(f(x)+g(x))=\frac{d}{d x} f(x)+\frac{d}{d x} g(x) .
\]

We now ask the same question for multiplication since we have the same limit law:
\[
\lim _{x \rightarrow a}(f(x) g(x))=\left(\lim _{x \rightarrow a} f(x)\right)\left(\lim _{x \rightarrow a} g(x)\right)
\]

It turns out the same property does \emph{not} hold for multiplication (and division)! In particular we have $(f g)^{\prime} \neq f^{\prime} g^{\prime}$ like we would want.

\begin{example}\label{example:5-8}
Let's look at a super simple example for proof. Let $f(x)=x$ and $g(x)=x$. Then,
% work: (fg)(x) = x^2 so (fg)'(x) = 2x, while f'(x) g'(x) = 1 * 1 = 1
\[
(fg)(x)=x\cdot x=x^{2}, \qquad (fg)^{\prime}(x)=\answer{2x}, \qquad f^{\prime}(x) g^{\prime}(x)=\answer{1}
\]
\end{example}

So we need to figure something else out to calculate $(f g)^{\prime}$. This is a bit complicated, but it turns out the best way is to use geometry.

First, let's look at one of our functions.
\begin{image}
\includegraphics[width=0.5\textwidth]{b48a5324-93dc-45ee-b7c2-dc0b5706bda6-065_711_741_1506_690}
\end{image}
\begin{image}

\includegraphics[width=0.4\textwidth]{b48a5324-93dc-45ee-b7c2-dc0b5706bda6-066_400_403_699_860}
\end{image}
% geometric derivation of the product rule done live in the original (white space); open work per spec 3(e)
\begin{freeResponse}\end{freeResponse}

This is called the \emph{product rule}: If $f$ and $g$ are both differentiable, then
\[
\frac{d}{dx}\left(f(x) g(x)\right)=\answer{f^{\prime}(x) g(x)+f(x) g^{\prime}(x)}
\]

\begin{example}\label{example:5-9}
Let's do an example. Let $f(x)=x^{2} e^{x}$. What is $f^{\prime}(x)$?
\[
f^{\prime}(x)=\answer{2 x e^{x}+x^{2} e^{x}}
\]
\end{example}

We can also do quotients based off a similar technique as the product.
% derivation of the quotient rule done live in the original (white space); open work per spec 3(e)
\begin{freeResponse}\end{freeResponse}

This is called the \emph{quotient rule}: If $f$ and $g$ are both differentiable and $g(x) \neq 0$, then
\[
\frac{d}{dx}\left(\frac{f(x)}{g(x)}\right)=\answer{\frac{f^{\prime}(x) g(x)-f(x) g^{\prime}(x)}{g(x)^{2}}}
\]

\begin{warning}
Note that in the quotient rule, the \emph{order} matters since we are subtracting. Make sure to keep this in mind.
\end{warning}

\begin{example}\label{example:5-10}
Let $f(x)=\frac{x^{3}+3 x}{x^{2}+1}$.
% work: numerator' = 3x^2+3, denominator' = 2x;
% f' = ((3x^2+3)(x^2+1) - (x^3+3x)(2x)) / (x^2+1)^2 = (x^4+3)/(x^2+1)^2
\[
f^{\prime}(x)=\frac{\answer{3 x^{2}+3}\cdot\left(x^{2}+1\right)-\left(x^{3}+3 x\right)\cdot\answer{2x}}{\answer{\left(x^{2}+1\right)^{2}}}
\]

Therefore,
\[
f^{\prime}(x)=\frac{x^{4}+3}{x^{4}+2 x^{2}+1}
\]
\end{example}

\begin{exercise}\label{exercise:5-11}
% AUTHOR CHOICE: example chosen during conversion, change freely
Now try this one: find $f^{\prime}(x)$ where $f(x)=\dfrac{e^{x}}{x^{2}+1}$.
% work: f' = (e^x (x^2+1) - e^x (2x)) / (x^2+1)^2 = e^x (x^2 - 2x + 1)/(x^2+1)^2 = e^x (x-1)^2/(x^2+1)^2
\[
f^{\prime}(x)=\answer{\frac{e^{x}(x-1)^{2}}{\left(x^{2}+1\right)^{2}}}
\]
\end{exercise}

\begin{example}\label{example:5-12}
Our next real world example comes from electrical engineering. There's a law in EE called ``Ohm's Law'' which states that the voltage is equal to the current times the resistance at any given point in a circuit,i.e., $V=I R$. We're going to try and calculate the instantaneous voltage at a given time if both the current and the resistance is changing over time.

If $I=i(t)$ is the function for the amount of current at any given time, then the average current over time is given by:
\[
\text{average current}=\frac{\Delta I}{\Delta t}
\]

and the instantaneous current at any given time is given by:
\[
\text{current}=\lim _{\Delta t \rightarrow 0} \frac{\Delta I}{\Delta t}=\frac{d I}{d t}=i^{\prime}(t)
\]

Similarly, if $R=r(t)$ is the function for the resistance at any given time, then the average resistance over time is given by:
\[
\text{average resistance}=\frac{\Delta R}{\Delta t}
\]

and the instantaneous resistance at any given time is given by:
\[
\text{resistance}=\lim _{\Delta t \rightarrow 0} \frac{\Delta R}{\Delta t}=\frac{d R}{d t}=r^{\prime}(t)
\]

We know that the average rate of voltage over time is given by:
\[
\text{average voltage}=\frac{\Delta V}{\Delta t}
\]

and, by what we saw earlier, the instantaneous voltage is given by:
\[
\text{voltage}=\lim _{\Delta t \rightarrow 0} \frac{\Delta V}{\Delta t}=\frac{d V}{d t}=\answer{i^{\prime}(t) r(t)+i(t) r^{\prime}(t)}
\]

Let's take as an example $I=i(t)=t^{2}$ and $R=r(t)=e^{t}$. Then what is the instantaneous voltage when $t=2$?
% work: V = t^2 e^t, dV/dt = 2t e^t + t^2 e^t; at t = 2: 4e^2 + 4e^2 = 8e^2
\[
\frac{d V}{d t}=\answer{2 t e^{t}+t^{2} e^{t}}, \qquad \left.\frac{d V}{d t}\right|_{t=2}=\answer{8 e^{2}}
\]
\end{example}

\end{document}
